\documentclass[12pt]{article}
\usepackage[utf8]{inputenc}
\usepackage[spanish]{babel}
\usepackage{amsmath}
\usepackage{geometry}
\geometry{margin=2.5cm}

\title{Guía de Ejercicios 07 – Notación Científica}
\author{Fundamentos de Matemática Básica \\ Agronomía}
\date{}

\begin{document}
	
	\maketitle
	
	\section*{I. Escribe en notación científica los siguientes valores relacionados con el área de la Agronomía.}
	
	\begin{enumerate}
		\item Número de granos cosechados en una temporada: 12.000.000.000
		\item Microorganismos presentes en una muestra de suelo fértil: 3.568.987.000.000
		\item Costo total de implementación de un sistema de riego tecnificado: 140.005.000.000
		\item Producción de semillas en una unidad experimental: 1.983
		\item Cantidad de plantas establecidas en un ensayo agrícola: 407.000
		\item Longitud promedio de una raíz principal (en m): 0,5480350
		\item Concentración de micronutriente en una solución fertilizante: 0,0002501
		\item Caudal de agua aplicado por goteo en una medición (en m$^3$): 0,00006
		\item Tamaño promedio de una partícula de arcilla en el suelo (en m): 0,0000000000045
		\item Precisión de lectura de un sensor de humedad: 150,000015
		\item Masa de una muestra de suelo analizada (en g): 125,05
		\item Distancia recorrida por una máquina agrícola en una jornada (en m): 3018,541
	\end{enumerate}
	
	\section*{II. Escriba el valor numérico correspondiente a los siguientes datos en notación científica.}
	
	\begin{enumerate}
		\item $2{,}85 \cdot 10^{12}$
		\item $1{,}0004 \cdot 10^6$
		\item $5 \cdot 10^{-8}$
		\item $8{,}205 \cdot 10^{-5}$
		\item $3{,}141529 \cdot 10^3$
		\item $3{,}1 \cdot 10^{-4}$
		\item $1 \cdot 10^{-10}$
		\item $3 \cdot 10^4$
		\item $4{,}30051 \cdot 10^{-2}$
		\item $9{,}09 \cdot 10^2$
		\item $4{,}03020109 \cdot 10^5$
		\item $6{,}7 \cdot 10^{-5}$
	\end{enumerate}
	
	\section*{III. Resuelve los siguientes ejercicios y expresa el resultado en notación científica.}
	
	\begin{enumerate}
		\item $2{,}54 \cdot 10^5 + 5{,}4 \cdot 10^5$
		\item $3 \cdot 10^6 - 2{,}4 \cdot 10^6$
		\item $1{,}15 \cdot 10^{15} - 4 \cdot 10^{15}$
		\item $6{,}05 \cdot 10^9 + 5 \cdot 10^8$
		\item $5{,}438 \cdot 10^{17} - 1{,}2 \cdot 10^{18}$
		\item $8 \cdot 10^5 - 3 \cdot 10^5 + 12 \cdot 10^4$
		\item $5{,}5 \cdot 10^{21} \cdot 1{,}8 \cdot 10^{15}$
		\item $6 \cdot 10^8 - 3 \cdot 10^5 + 12 \cdot 10^4$
		\item $4{,}5 \cdot 10^{11} \cdot 4 \cdot 10^{19}$
		\item $\frac{3{,}6 \cdot 10^{16}}{6 \cdot 10^{14}}$
		\item $\frac{2{,}43 \cdot 10^{54}}{2{,}7 \cdot 10^{45}}$
		\item $\frac{6{,}4 \cdot 10^{102}}{2 \cdot 10^{85}}$
	\end{enumerate}
	
	\section*{IV. Problemas Aplicados a Agronomía}
	
	\begin{enumerate}
		\item Un agrónomo debe aplicar fertilizante en una parcela con una dosis de $2{,}5 \cdot 10^3$ gramos por hectárea. El terreno recibe 3 aplicaciones por semana durante 2 semanas. Además, antes del plan principal se realizó una aplicación inicial de $1 \cdot 10^3$ gramos. \\
		¿Cuál es la cantidad total de fertilizante aplicada?
		
		\item Un cultivo presenta el siguiente balance hídrico diario: \\
		Agua aportada por riego y lluvia: $3{,}2 \cdot 10^3$ litros \\
		Pérdida por evaporación y drenaje: $1{,}8 \cdot 10^3$ litros \\
		Consumo por absorción de las plantas: $2{,}5 \cdot 10^3$ litros \\
		¿Cuál es la ganancia o pérdida neta de agua en el sistema?
		
		\item Durante un monitoreo de semillas se registran dos lotes: \\
		Primer lote: densidad de $3{,}0 \cdot 10^4$ semillas por bandeja, durante un área sembrada de $2{,}0 \cdot 10^{-2}$ hectáreas \\
		Segundo lote: densidad de $2{,}0 \cdot 10^4$ semillas por bandeja, en la misma superficie \\
		\begin{itemize}
			\item[a)] ¿Cuántas semillas corresponden a cada lote?
			\item[b)] ¿Cuál es la cantidad total de semillas utilizadas?
		\end{itemize}
		
		\item Un agrónomo administra un producto fitosanitario a dos cultivos experimentales: \\
		Parcela de maíz: $4{,}31 \cdot 10^2$ m$^2$ \\
		Parcela de trigo: $1{,}24 \cdot 10^2$ m$^2$ \\
		Dosis: $2{,}0 \cdot 10^{-1}$ litros/m$^2$ \\
		\begin{itemize}
			\item[a)] ¿Cuántos litros recibe cada parcela?
			\item[b)] ¿Cuánta más dosis recibió la parcela de maíz?
		\end{itemize}
	\end{enumerate}
	
\end{document}